\section*{Capítulo \ref{ch:g8:the-atmosphere} --- A atmosfera}

\begin{solution}{exo:g8:the-atmosphere:1}
(a) \omterm{def:g8:the-atmosphere:layers}{estratosfera}; (b) \omterm{def:g8:the-atmosphere:layers}{troposfera}; (c) \omterm{def:g8:the-atmosphere:layers}{mesosfera}; (d) \omterm{def:g8:the-atmosphere:layers}{termosfera}.
\end{solution}

\begin{solution}{exo:g8:the-atmosphere:2}
(a) \qty{15}{\degreeCelsius}. (b) Cerca de \qty{-56}{\degreeCelsius}.
(c) Cerca de \qty{-3}{\degreeCelsius}. (d) Cerca de \qty{-86}{\degreeCelsius}.
(e) De \qty{11}{km} até cerca de \qty{20}{km} (e ao longo de alguns
quilômetros perto de \qty{50}{km} e logo acima de \qty{86}{km}).
\end{solution}

\begin{solution}{exo:g8:the-atmosphere:3}
(a) Falso: ela cai na \omterm{def:g8:the-atmosphere:layers}{troposfera}, sobe na \omterm{def:g8:the-atmosphere:layers}{estratosfera}, cai na \omterm{def:g8:the-atmosphere:layers}{mesosfera}
e sobe de novo na \omterm{def:g8:the-atmosphere:layers}{termosfera}. (b) Verdadeiro. (c) Falso: ela fica na
\omterm{def:g8:the-atmosphere:layers}{estratosfera}. (d) Verdadeiro.
\end{solution}

\begin{solution}{exo:g8:the-atmosphere:4}
(a) Na \omterm{def:g8:the-atmosphere:layers}{troposfera}. (b) Na \omterm{def:g8:the-atmosphere:layers}{estratosfera} (na atmosfera padrão, a \omterm{def:g8:the-atmosphere:layers}{tropopausa}
fica a \qty{11}{km}). (c) Na \omterm{def:g8:the-atmosphere:layers}{estratosfera}. (d) Na \omterm{def:g8:the-atmosphere:layers}{mesosfera}.
\end{solution}

\begin{solution}{exo:g8:the-atmosphere:5}
(a) $20 - 6.5 \times 2 = 20 - 13 = \qty{7}{\degreeCelsius}$. (b) $20 -
6.5 \times 3 = 20 - 19.5 = \qty{0.5}{\degreeCelsius}$. (c) O ar precisa
esfriar 20 graus: $20 \div 6.5 \approx 3.1$, então a cerca de
\qty{3.1}{km}.
\end{solution}

\begin{solution}{exo:g8:the-atmosphere:6}
(a) $3 - 0.5 = \qty{2.5}{km}$. (b) $6.5 \times 2.5 = 16.25$: cerca de 16
graus. (c) $18 - 16.25 \approx \qty{2}{\degreeCelsius}$.
\end{solution}

\begin{solution}{exo:g8:the-atmosphere:7}
(a) \qty{500}{hPa}, \qty{250}{hPa} e \qty{125}{hPa}. (b) Perto do chão, a
regra é boa (\qty{500}{hPa} contra \qty{506}{hPa}). Mais acima, ela dá
demais: \qty{250}{hPa} contra \qty{227}{hPa}, e \qty{125}{hPa} contra
\qty{96}{hPa}. A regra da metade é só uma indicação aproximada.
\end{solution}

\begin{solution}{exo:g8:the-atmosphere:8}
(a) $10~\text{hPa} = 10^1$~hPa; $10^3 \div 10^1 = 10^2$: cem vezes
menor. (b) $0.01~\text{hPa} = 10^{-2}$~hPa. (c) $10^3 \div 10^{-2} =
10^5$: cem mil vezes menor.
\end{solution}

\begin{solution}{exo:g8:the-atmosphere:9}
(a) $6 \div 1.5 = 4$ vezes mais largo. (b) Enquanto ele sobe, o ar em
volta faz cada vez menos pressão, então o gás de dentro se espalha e o
balão cresce até a sua pele se romper. (c) Depois que o balão estoura, o
paraquedas freia a queda da radiossonda, para que ela não machuque
ninguém nem danifique nada.
\end{solution}

\begin{solution}{exo:g8:the-atmosphere:10}
(a) $20 \div 7 \approx 2.9$ vezes mais espessa. (b) Mais perto dos \omterm{def:g3:continents-and-oceans:poles}{polos}:
\qty{11}{km} fica \qty{4}{km} acima de \qty{7}{km}, mas \qty{9}{km} abaixo de \qty{20}{km} (ela descreve as \omterm{def:g6:latitude-longitude:latitude}{latitudes} médias, e não o \omterm{def:g3:continents-and-oceans:equator}{equador}).
(c) No começo, a \omterm{def:g8:the-atmosphere:layers}{tropopausa} fica a cerca de \qty{7}{km}, abaixo do avião:
o avião está na \omterm{def:g8:the-atmosphere:layers}{estratosfera}. No fim, a \omterm{def:g8:the-atmosphere:layers}{tropopausa} fica a cerca de
\qty{20}{km}: o avião está na \omterm{def:g8:the-atmosphere:layers}{troposfera}.
\end{solution}

\begin{solution}{exo:g8:the-atmosphere:11}
(a) $1 - 506 \div 1013 \approx 0.50$: \textbf{50~\%}. (b) $1 - 227 \div
1013 \approx 0.78$: \textbf{78~\%}. (c) $1 - 12 \div 1013 \approx 0.99$:
\textbf{99~\%}. (d) Cerca de 1~\%.
\end{solution}

\begin{solution}{exo:g8:the-atmosphere:12}
(a) $6371 \div 15 \approx 425$~km para cada centímetro. (b) $100 \div 425
\approx 0.24$~cm, cerca de \qty{2.4}{mm}. (c) $11 \div 425 \approx
0.026$~cm, cerca de \qty{0.26}{mm}. (d) $0.26 \div 0.1 \approx 2.6$: duas
ou três folhas de papel.
\end{solution}

\begin{solution}{pb:g8:the-atmosphere:1}
\textbf{1.} $300 \times 10 = 3000$~m.

\textbf{2.} $6000 \div 300 = 20$ minutos; $\num{30000} \div 300 = 100$
minutos.

\textbf{3.} 100 minutos, ou seja, 1~hora e 40~minutos.

\textbf{4.} A \qty{8}{km}: a \omterm{def:g8:the-atmosphere:layers}{troposfera} (a temperatura ainda está
caindo). A \qty{30}{km}: a \omterm{def:g8:the-atmosphere:layers}{estratosfera}.

\textbf{5.} $15 - (-50) = 65$ graus.

\textbf{6.} $65 \div 10 = \qty{6.5}{\degreeCelsius}$ por quilômetro.

\textbf{7.} De \qty{12}{km} a \qty{20}{km}, ela fica em
\qty{-56}{\degreeCelsius}. O resfriamento para entre \qty{10}{km} e
\qty{12}{km}: a \omterm{def:g8:the-atmosphere:layers}{tropopausa} fica a cerca de \qty{11}{km}.

\textbf{8.} Ela sobe: \qty{-55}{\degreeCelsius} a \qty{22}{km},
\qty{-47}{\degreeCelsius} a \qty{30}{km}. O balão está na \omterm{def:g8:the-atmosphere:layers}{estratosfera},
aquecida pelo ozônio que absorve a luz ultravioleta do Sol.

\textbf{9.} $1013 \div 2 \approx \qty{507}{hPa}$; ela é atingida entre
\qty{4}{km} (\qty{617}{hPa}) e \qty{6}{km} (\qty{472}{hPa}).

\textbf{10.} $1013 \div 4 \approx \qty{253}{hPa}$.

\textbf{11.} Entre \qty{10}{km} (\qty{265}{hPa}) e \qty{12}{km}
(\qty{194}{hPa}).

\textbf{12.} A partir de \qty{10}{km}, a pressão precisa cair mais $265 - 253 =
12$~hPa, dos 71~hPa que ela perde em 2~km: $12 \div 71 \times 2
\approx 0.34$~km. A pressão é um quarto do seu valor no \omterm{def:g4:world-physical-map:altitude}{nível do mar} a
cerca de \textbf{\qty{10.3}{km}}.
\end{solution}
